%% kapa: exam=1 round=34 year=2023 period=1 form=A subject=civil_law subject_ko=민법 seq=7 number=7
%% source: Q-Net 감정평가사 기출문제 (cst00309) artlSeq=5211789 file=raw/1008/5211789/2023년도 제34회 감정평가사 1차시험 1교시.pdf
%% conversion: kordoc 4.24.51 -> tools/gen1.py (auto); review=proofread-sonnet; flags=none
\begin{question}{7}
민법상 법인의 기관에 관한 설명으로 옳은 것은? (다툼이 있으면 판례에 따름)
\begin{choices}
  \item 이사의 변경등기는 대항요건이 아니라 효력발생요건이다.
  \item 이사가 수인인 경우, 특별한 사정이 없는 한 법인의 사무에 관하여 이사는 공동으로 법인을 대표한다.
  \item 사단법인의 정관 변경에 관한 사원총회의 권한은 정관에 의해 박탈할 수 있다.
  \item 이사회에서 법인과 어느 이사와의 관계사항을 의결하는 경우, 그 이사는 의사정족수 산정의 기초가 되는 이사의 수에 포함된다.
  \item 법인의 대표권 제한에 관한 사항이 등기되지 않았더라도 법인은 대표권 제한에 대해 악의인 제3자에게 대항할 수 있다.
\end{choices}
\end{question}
\answer{4}
